% ============================================================
%  CHAPITRE DE DÉMONSTRATION
%
%  Ce fichier n'est pas inclus par défaut.
%  Pour l'afficher : décommentez % ============================================================
%  CHAPITRE DE DÉMONSTRATION
%
%  Ce fichier n'est pas inclus par défaut.
%  Pour l'afficher : décommentez % ============================================================
%  CHAPITRE DE DÉMONSTRATION
%
%  Ce fichier n'est pas inclus par défaut.
%  Pour l'afficher : décommentez % ============================================================
%  CHAPITRE DE DÉMONSTRATION
%
%  Ce fichier n'est pas inclus par défaut.
%  Pour l'afficher : décommentez \input{chapitres/demo} dans main.tex.
%  Il montre tous les éléments disponibles dans le template.
% ============================================================

\chapitre{3}{Démonstration : toutes les fonctionnalités}

\begin{objectifs}
\begin{itemize}
  \item Voir un exemple de chaque élément disponible.
  \item Savoir quel environnement utiliser pour chaque besoin.
\end{itemize}
\end{objectifs}

\soussection{3.1}{Texte, listes et mise en forme}

Un paragraphe normal s'écrit directement. On peut mettre un mot en
\textbf{gras}, un autre en \textit{italique}, ou afficher du
\code{code inline}. Le terme \motcle{important} est coloré.

\begin{itemize}
  \item Premier point.
  \item Deuxième point avec une liste imbriquée :
  \begin{enumerate}
    \item Sous-point un.
    \item Sous-point deux.
  \end{enumerate}
\end{itemize}

\begin{definition}
Une \textbf{définition} présente un concept nouveau de manière rigoureuse
et précise.
\end{definition}

\begin{retenir}
Ce bloc sert à résumer ce qu'il faut absolument retenir.
\end{retenir}

\begin{attention}
Ce bloc signale un point dangereux ou une erreur à ne pas commettre.
\end{attention}

\begin{piege}
Ce bloc décrit un piège classique et la manière de l'éviter.
\end{piege}

\begin{conseil}
Ce bloc donne une bonne pratique ou une astuce.
\end{conseil}

\begin{choix}
Ce bloc justifie un choix technique recommandé.
\end{choix}

\begin{atelier}
Ce bloc propose une activité pratique ou un exercice guidé.
\end{atelier}

\soussection{3.2}{Blocs de code}

Voici un exemple Java :

\begin{lstlisting}[style=java,caption={Exemple de contrôleur Java}]
@RestController
@RequestMapping("/api/articles")
public class ArticleController {

    @GetMapping("/{id}")
    public ResponseEntity<Article> findById(@PathVariable Long id) {
        return ResponseEntity.ok(service.findById(id));
    }
}
\end{lstlisting}

Un exemple Python :

\begin{lstlisting}[style=python,caption={Exemple Python}]
def salutation(nom: str) -> str:
    return f"Bonjour, {nom} !"

print(salutation("Alice"))
\end{lstlisting}

Un exemple JSON :

\begin{lstlisting}[style=json,caption={Exemple de réponse JSON}]
{
  "id": 42,
  "titre": "Mon article",
  "publie": true
}
\end{lstlisting}

Une commande shell :

\begin{lstlisting}[style=bash,caption={Commandes de compilation}]
lualatex main.tex
lualatex main.tex
\end{lstlisting}

Un fichier de configuration :

\begin{lstlisting}[style=properties,caption={Fichier de configuration}]
serveur.port=8080
serveur.hote=localhost
mode.debug=true
\end{lstlisting}

\soussection{3.3}{Tableaux et mathématiques}

\begin{table}[htbp]
\centering
\small
\begin{tabularx}{\textwidth}{lX}
\toprule
\textbf{Élément} & \textbf{Description} \\
\midrule
\code{lstlisting} & Affiche un bloc de code. \\
\code{tcolorbox} & Affiche un encadré coloré. \\
\code{tabularx} & Affiche un tableau qui s'adapte à la largeur. \\
\bottomrule
\end{tabularx}
\caption{Exemple de tableau avec booktabs et tabularx}
\end{table}

Une formule mathématique :

\begin{equation*}
  f(x) = \sum_{i=1}^{n} x_i^2
\end{equation*}

Un lien : \href{https://example.com}{https://example.com}.

\soussection{3.4}{Conclusion}

\begin{resume}
Ce chapitre de démonstration montre tous les éléments de base du
template. Il n'est pas destiné à être conservé dans le cours final.
\end{resume}
 dans main.tex.
%  Il montre tous les éléments disponibles dans le template.
% ============================================================

\chapitre{3}{Démonstration : toutes les fonctionnalités}

\begin{objectifs}
\begin{itemize}
  \item Voir un exemple de chaque élément disponible.
  \item Savoir quel environnement utiliser pour chaque besoin.
\end{itemize}
\end{objectifs}

\soussection{3.1}{Texte, listes et mise en forme}

Un paragraphe normal s'écrit directement. On peut mettre un mot en
\textbf{gras}, un autre en \textit{italique}, ou afficher du
\code{code inline}. Le terme \motcle{important} est coloré.

\begin{itemize}
  \item Premier point.
  \item Deuxième point avec une liste imbriquée :
  \begin{enumerate}
    \item Sous-point un.
    \item Sous-point deux.
  \end{enumerate}
\end{itemize}

\begin{definition}
Une \textbf{définition} présente un concept nouveau de manière rigoureuse
et précise.
\end{definition}

\begin{retenir}
Ce bloc sert à résumer ce qu'il faut absolument retenir.
\end{retenir}

\begin{attention}
Ce bloc signale un point dangereux ou une erreur à ne pas commettre.
\end{attention}

\begin{piege}
Ce bloc décrit un piège classique et la manière de l'éviter.
\end{piege}

\begin{conseil}
Ce bloc donne une bonne pratique ou une astuce.
\end{conseil}

\begin{choix}
Ce bloc justifie un choix technique recommandé.
\end{choix}

\begin{atelier}
Ce bloc propose une activité pratique ou un exercice guidé.
\end{atelier}

\soussection{3.2}{Blocs de code}

Voici un exemple Java :

\begin{lstlisting}[style=java,caption={Exemple de contrôleur Java}]
@RestController
@RequestMapping("/api/articles")
public class ArticleController {

    @GetMapping("/{id}")
    public ResponseEntity<Article> findById(@PathVariable Long id) {
        return ResponseEntity.ok(service.findById(id));
    }
}
\end{lstlisting}

Un exemple Python :

\begin{lstlisting}[style=python,caption={Exemple Python}]
def salutation(nom: str) -> str:
    return f"Bonjour, {nom} !"

print(salutation("Alice"))
\end{lstlisting}

Un exemple JSON :

\begin{lstlisting}[style=json,caption={Exemple de réponse JSON}]
{
  "id": 42,
  "titre": "Mon article",
  "publie": true
}
\end{lstlisting}

Une commande shell :

\begin{lstlisting}[style=bash,caption={Commandes de compilation}]
lualatex main.tex
lualatex main.tex
\end{lstlisting}

Un fichier de configuration :

\begin{lstlisting}[style=properties,caption={Fichier de configuration}]
serveur.port=8080
serveur.hote=localhost
mode.debug=true
\end{lstlisting}

\soussection{3.3}{Tableaux et mathématiques}

\begin{table}[htbp]
\centering
\small
\begin{tabularx}{\textwidth}{lX}
\toprule
\textbf{Élément} & \textbf{Description} \\
\midrule
\code{lstlisting} & Affiche un bloc de code. \\
\code{tcolorbox} & Affiche un encadré coloré. \\
\code{tabularx} & Affiche un tableau qui s'adapte à la largeur. \\
\bottomrule
\end{tabularx}
\caption{Exemple de tableau avec booktabs et tabularx}
\end{table}

Une formule mathématique :

\begin{equation*}
  f(x) = \sum_{i=1}^{n} x_i^2
\end{equation*}

Un lien : \href{https://example.com}{https://example.com}.

\soussection{3.4}{Conclusion}

\begin{resume}
Ce chapitre de démonstration montre tous les éléments de base du
template. Il n'est pas destiné à être conservé dans le cours final.
\end{resume}
 dans main.tex.
%  Il montre tous les éléments disponibles dans le template.
% ============================================================

\chapitre{3}{Démonstration : toutes les fonctionnalités}

\begin{objectifs}
\begin{itemize}
  \item Voir un exemple de chaque élément disponible.
  \item Savoir quel environnement utiliser pour chaque besoin.
\end{itemize}
\end{objectifs}

\soussection{3.1}{Texte, listes et mise en forme}

Un paragraphe normal s'écrit directement. On peut mettre un mot en
\textbf{gras}, un autre en \textit{italique}, ou afficher du
\code{code inline}. Le terme \motcle{important} est coloré.

\begin{itemize}
  \item Premier point.
  \item Deuxième point avec une liste imbriquée :
  \begin{enumerate}
    \item Sous-point un.
    \item Sous-point deux.
  \end{enumerate}
\end{itemize}

\begin{definition}
Une \textbf{définition} présente un concept nouveau de manière rigoureuse
et précise.
\end{definition}

\begin{retenir}
Ce bloc sert à résumer ce qu'il faut absolument retenir.
\end{retenir}

\begin{attention}
Ce bloc signale un point dangereux ou une erreur à ne pas commettre.
\end{attention}

\begin{piege}
Ce bloc décrit un piège classique et la manière de l'éviter.
\end{piege}

\begin{conseil}
Ce bloc donne une bonne pratique ou une astuce.
\end{conseil}

\begin{choix}
Ce bloc justifie un choix technique recommandé.
\end{choix}

\begin{atelier}
Ce bloc propose une activité pratique ou un exercice guidé.
\end{atelier}

\soussection{3.2}{Blocs de code}

Voici un exemple Java :

\begin{lstlisting}[style=java,caption={Exemple de contrôleur Java}]
@RestController
@RequestMapping("/api/articles")
public class ArticleController {

    @GetMapping("/{id}")
    public ResponseEntity<Article> findById(@PathVariable Long id) {
        return ResponseEntity.ok(service.findById(id));
    }
}
\end{lstlisting}

Un exemple Python :

\begin{lstlisting}[style=python,caption={Exemple Python}]
def salutation(nom: str) -> str:
    return f"Bonjour, {nom} !"

print(salutation("Alice"))
\end{lstlisting}

Un exemple JSON :

\begin{lstlisting}[style=json,caption={Exemple de réponse JSON}]
{
  "id": 42,
  "titre": "Mon article",
  "publie": true
}
\end{lstlisting}

Une commande shell :

\begin{lstlisting}[style=bash,caption={Commandes de compilation}]
lualatex main.tex
lualatex main.tex
\end{lstlisting}

Un fichier de configuration :

\begin{lstlisting}[style=properties,caption={Fichier de configuration}]
serveur.port=8080
serveur.hote=localhost
mode.debug=true
\end{lstlisting}

\soussection{3.3}{Tableaux et mathématiques}

\begin{table}[htbp]
\centering
\small
\begin{tabularx}{\textwidth}{lX}
\toprule
\textbf{Élément} & \textbf{Description} \\
\midrule
\code{lstlisting} & Affiche un bloc de code. \\
\code{tcolorbox} & Affiche un encadré coloré. \\
\code{tabularx} & Affiche un tableau qui s'adapte à la largeur. \\
\bottomrule
\end{tabularx}
\caption{Exemple de tableau avec booktabs et tabularx}
\end{table}

Une formule mathématique :

\begin{equation*}
  f(x) = \sum_{i=1}^{n} x_i^2
\end{equation*}

Un lien : \href{https://example.com}{https://example.com}.

\soussection{3.4}{Conclusion}

\begin{resume}
Ce chapitre de démonstration montre tous les éléments de base du
template. Il n'est pas destiné à être conservé dans le cours final.
\end{resume}
 dans main.tex.
%  Il montre tous les éléments disponibles dans le template.
% ============================================================

\chapitre{3}{Démonstration : toutes les fonctionnalités}

\begin{objectifs}
\begin{itemize}
  \item Voir un exemple de chaque élément disponible.
  \item Savoir quel environnement utiliser pour chaque besoin.
\end{itemize}
\end{objectifs}

\soussection{3.1}{Texte, listes et mise en forme}

Un paragraphe normal s'écrit directement. On peut mettre un mot en
\textbf{gras}, un autre en \textit{italique}, ou afficher du
\code{code inline}. Le terme \motcle{important} est coloré.

\begin{itemize}
  \item Premier point.
  \item Deuxième point avec une liste imbriquée :
  \begin{enumerate}
    \item Sous-point un.
    \item Sous-point deux.
  \end{enumerate}
\end{itemize}

\begin{definition}
Une \textbf{définition} présente un concept nouveau de manière rigoureuse
et précise.
\end{definition}

\begin{retenir}
Ce bloc sert à résumer ce qu'il faut absolument retenir.
\end{retenir}

\begin{attention}
Ce bloc signale un point dangereux ou une erreur à ne pas commettre.
\end{attention}

\begin{piege}
Ce bloc décrit un piège classique et la manière de l'éviter.
\end{piege}

\begin{conseil}
Ce bloc donne une bonne pratique ou une astuce.
\end{conseil}

\begin{choix}
Ce bloc justifie un choix technique recommandé.
\end{choix}

\begin{atelier}
Ce bloc propose une activité pratique ou un exercice guidé.
\end{atelier}

\soussection{3.2}{Blocs de code}

Voici un exemple Java :

\begin{lstlisting}[style=java,caption={Exemple de contrôleur Java}]
@RestController
@RequestMapping("/api/articles")
public class ArticleController {

    @GetMapping("/{id}")
    public ResponseEntity<Article> findById(@PathVariable Long id) {
        return ResponseEntity.ok(service.findById(id));
    }
}
\end{lstlisting}

Un exemple Python :

\begin{lstlisting}[style=python,caption={Exemple Python}]
def salutation(nom: str) -> str:
    return f"Bonjour, {nom} !"

print(salutation("Alice"))
\end{lstlisting}

Un exemple JSON :

\begin{lstlisting}[style=json,caption={Exemple de réponse JSON}]
{
  "id": 42,
  "titre": "Mon article",
  "publie": true
}
\end{lstlisting}

Une commande shell :

\begin{lstlisting}[style=bash,caption={Commandes de compilation}]
lualatex main.tex
lualatex main.tex
\end{lstlisting}

Un fichier de configuration :

\begin{lstlisting}[style=properties,caption={Fichier de configuration}]
serveur.port=8080
serveur.hote=localhost
mode.debug=true
\end{lstlisting}

\soussection{3.3}{Tableaux et mathématiques}

\begin{table}[htbp]
\centering
\small
\begin{tabularx}{\textwidth}{lX}
\toprule
\textbf{Élément} & \textbf{Description} \\
\midrule
\code{lstlisting} & Affiche un bloc de code. \\
\code{tcolorbox} & Affiche un encadré coloré. \\
\code{tabularx} & Affiche un tableau qui s'adapte à la largeur. \\
\bottomrule
\end{tabularx}
\caption{Exemple de tableau avec booktabs et tabularx}
\end{table}

Une formule mathématique :

\begin{equation*}
  f(x) = \sum_{i=1}^{n} x_i^2
\end{equation*}

Un lien : \href{https://example.com}{https://example.com}.

\soussection{3.4}{Conclusion}

\begin{resume}
Ce chapitre de démonstration montre tous les éléments de base du
template. Il n'est pas destiné à être conservé dans le cours final.
\end{resume}
